% ==============================================================
% SECTION VIII: CONCLUSION
% ==============================================================
\section{Conclusion}

We presented NeuroAgent, a three-phase, eight-agent system that unifies the brain tumor diagnostic chain---from raw multi-modal MRI to clinician-ready structured reports---within a single LangGraph-orchestrated pipeline with persistent clinical memory. Evaluated on 369 BraTS 2020 cases (all metrics sourced from \texttt{brats\_outputs/evaluation\_results.csv}), the system achieves mean DSC $0.8517 \pm 0.1055$, median HD95 3.61~mm, mean sensitivity $0.9384$, and 99.7\% pipeline completion at $3.65 \pm 0.16$~seconds per case. The WHO CNS5 classifier produces morphologically consistent classifications with sigmoid-calibrated confidence (mean 0.902) across five tumor types, while the RANO assessment agent correctly defaults to conservative outputs when longitudinal data is unavailable.

Three architectural principles distinguish NeuroAgent from prior systems. First, clinical protocols (WHO CNS5, RANO, CAP) are encoded as deterministic agent logic rather than delegated to LLM interpretation, eliminating hallucination risk in protocol-critical outputs. Second, every agent implements graceful two-tier fallback, enabling deployment without GPU, Ollama, or Weaviate dependencies while preserving clinical content. Third, the ChromaDB-backed patient memory accumulates per-patient scan history, classification trajectories, and physician corrections, enabling progression-aware reasoning that no prior segmentation or classification system provides.

Future work will pursue four directions: (i) multi-class subregion segmentation evaluating ET, TC, and WT independently to provide more granular surgical planning information; (ii) prospective clinical validation against biopsy-confirmed histopathological diagnoses to quantify WHO classification accuracy; (iii) multi-timepoint evaluation on longitudinal datasets to demonstrate the patient memory system's progression tracking capabilities; and (iv) federated memory architectures that enable cross-institutional knowledge sharing while preserving patient privacy through differential privacy mechanisms.
